\chapter{Leadership Roles}\label{ch:in:leadership-roles}

In the United Kingdom, the most senior people at a regulated firm hold senior management functions, need the regulator's
approval before they start, and each has a written statement of their responsibilities. The regulator's guide explains
why: they ``are the most senior people in a firm with the greatest potential to cause harm or impact upon market
integrity''. At the top of a trading firm, a job title is also a legal accountability. This chapter describes the
leadership roles a quantitative career can reach (\omterm{def:in:leadership-roles:head}{head of desk}, the chief officers, partner, chief executive and
founder), the accountability regimes that attach to them, and the economics of making partner.

\begin{center}\footnotesize
\begin{tabular}{@{}p{2.2cm}p{2.5cm}p{2.5cm}p{2.5cm}p{2.5cm}@{}}
\toprule
\multicolumn{5}{@{}l}{\textbf{Role cards: leadership}}\\
\midrule
 & \omterm{def:in:leadership-roles:head}{head of desk} & chief risk officer & partner & chief executive\\
\midrule
answers for & a desk's P\&L, risk and people & the firm's risk framework & the firm, as an owner & the whole business\\
paid by & salary and a bonus on the desk & salary and a control bonus & a share of profit by points & salary, bonus and ownership\\
UK function & (certified, not SMF) & SMF4 (enhanced firms) & SMF27 & SMF1\\
taught in & Book~16, ch.~6 & Book~16, ch.~12 & Book~16, ch.~2 & Book~16, ch.~1\\
\bottomrule
\end{tabular}
\end{center}

\section{Head of desk}

\begin{definition}[Head of desk]\label{def:in:leadership-roles:head}
A \emph{head of desk}\index{head of desk} is the person who runs one trading desk: sets its strategy within the
firm's risk appetite, allocates its limits and capital among its traders, answers for its profit and loss, its risk
and its people, and agrees its revenue budget with the business's management.
\end{definition}

The \omterm{def:in:leadership-roles:head}{head of desk} is the first leadership role most traders and \omterm{def:in:quant-researcher:def}{quantitative researchers} reach. The revenue budget and the
risk appetite of Book~16, chapter~6, are the terms of the job; hiring, pay recommendations and the desk's controls are
its content. A \omterm{def:in:leadership-roles:head}{head of desk} still trades at many firms, but the job's measure changes from one's own result to the
desk's. In the UK scheme, a \omterm{def:in:leadership-roles:head}{head of desk} who holds no senior management function is not a senior manager, and may fall
under the \omterm{def:in:trader:certification}{certification regime} of chapter~16 instead.

\section{Chief risk, technology and operating officers}

The chief officers run functions rather than books. The chief risk officer (Book~16, chapter~12) owns the risk
framework and reports to the board as well as the chief executive; the chief technology officer owns the systems that
trade (chapters~19 and~20); the chief operating officer owns operations and often finance and the firm's infrastructure
(chapter~24). Each can be the destination of a quantitative career that has broadened: a \omterm{def:in:bank-quant:risk}{risk quant} can become head of
market risk and then chief risk officer; an engineer, chief technology officer. In enhanced UK firms each of
these roles is a named senior management function: SMF4 for the chief risk officer, SMF24 for the chief operations
function, SMF2 for the chief finance function.

\section{Partner}

Some proprietary trading firms are partnerships (Book~16, chapter~2): a limited liability partnership whose members own
the firm, contribute capital, share its profit by points and leave part of it in the firm. Chapter~11 read one such
firm's accounts, where \omterm{def:in:reading-a-trading-firms-accounts:members}{members' remuneration} is a share of profit rather than a salary. Making partner changes three
things: pay becomes a share of the firm's profit rather than of one's own; part of it stays in the firm as capital at
risk; and the partner answers for the whole firm.

\begin{method}[A new partner's first years]\label{met:in:leadership-roles:partner}
Let the firm's profit in year $t$ be $\Pi_t$, and let the new partner hold $p$ of the partnership's $P_t$ points, with
$P_t$ rising when others are admitted. The partner's share is $s_t=p/P_t$; the partner draws $\rho\,s_t\max(\Pi_t,0)$ and
retains the rest of a positive share in their capital account; a loss reduces the capital account by $s_t\Pi_t$. Starting
capital $C_0$ is the contribution on admission.
\end{method}

\begin{example}[Making partner]\label{ex:in:leadership-roles:partner}
Illustrative terms: firm profit normal with mean \$300 million and standard deviation \$200 million a year; 30 partners
with 100 points each; the new partner receives 30 points and contributes \$1 million; 60\% of each positive share is
drawn; three more partners of 30 points are admitted in year three. Over five years the new partner's expected draws are
\$8.93 million (10th--90th percentile \$5.79--12.14 million), and their capital account is expected to reach \$6.67
million, with a 5th percentile of \$3.32 million. In 9.4\% of paths the capital account falls below the contribution at
some year-end. Against a salaried \omterm{def:in:leadership-roles:head}{head of desk} at \$1.2 million a year, the draws alone break even at about 20 points,
0.66\% of the partnership (\cref{fig:in:leadership-roles:partner}).
\end{example}

\begin{omfigure}
\begin{tikzpicture}
\begin{axis}[ybar,area legend,width=10.5cm,height=5.4cm,xmin=0.4,xmax=5.6,ymin=0,ymax=8,bar width=13pt,xtick={1,2,3,4,5},
  xlabel={\small year after admission},ylabel={\small \$ million},tick label style={font=\footnotesize},
  grid=major,grid style={gray!20},table/col sep=comma,
  legend style={font=\scriptsize,at={(0.03,0.97)},anchor=north west,legend cell align=left}]
\addplot[fill=omDef,draw=none] table[x=year,y=draw]{figdata/industry/25-leadership-roles/partner.csv};
\addplot[fill=omThm!60,draw=none] table[x=year,y=capital]{figdata/industry/25-leadership-roles/partner.csv};
\addplot[black,thick,sharp plot,mark=*,mark size=1.5pt,forget plot] table[x=year,y=cap_p05]{figdata/industry/25-leadership-roles/partner.csv};
\addplot[black!60,thick,dashed,sharp plot,domain=0.4:5.6,
  legend image code/.code={\draw[black!60,thick,dashed] (0cm,0cm) -- (0.6cm,0cm);}] {1.2};
\legend{expected draw,expected capital account,salaried alternative a year}
\end{axis}
\end{tikzpicture}

\omcaption{The illustrative new partner's expected draw each year and expected capital account at each year-end (the dots
are the capital account's 5th percentile), against a salaried \omterm{def:in:leadership-roles:head}{head of desk}'s \$1.2 million a year. The capital account
is the partner's money at risk in the firm. 20\,000 paths. Data: \texttt{in\_leader.summary}, from
\texttt{firm.roles.partner\_path} on \texttt{firm.partnership}.}\label{fig:in:leadership-roles:partner}
\end{omfigure}

The comparison needs care. The draws are income; the capital account is wealth that the partner will receive back only
on leaving, in instalments (Book~16, chapter~2), and can lose. A partner who counts the capital as pay is taking the
firm's risk with a large part of their savings, which is the partnership's purpose: the owners' money is the first to
absorb a loss.

\section{Chief executive and founder}

At the top the roles can merge: a founder may be chief executive, chief investment officer and the largest partner or
shareholder at once. Book~16, chapter~1, sets out the economics of the trading firm the founder builds, and
chapter~25 of that book the launch of a fund. For a quantitative career the route there is rarely through titles: it
runs through a strategy or a technology that works and the capital and people to scale it. The survey's chief executives
in the securities industry earned a median wage of \$431\,020 in May 2025, with a 90th percentile of \$749\,840; for
founders and partners wages are a small part of pay, which the survey does not measure.

\begin{dated}{2025-05}{What leaders are paid: the public evidence}\label{dat:in:leadership-roles:pay}
Occupational survey, May 2025, securities industry: chief executives 4\,700 employed, median wage \$431\,020 (10th--90th
percentile \$186\,460--749\,840); computer and information systems managers 18\,630, median \$214\,460; general and
operations managers 39\,880, median \$198\,900; financial managers median \$223\,860. EBA \omterm{def:in:pay-levels-by-role-firm-type-and-seniority:median}{high earners} 2024: 586 members
of credit institutions' management bodies (management function) paid over one million euros, average 2.10 million,
variable 0.73 of fixed; 61 at investment firms, average 1.99 million, variable 2.07 of fixed. Wages and aggregates only;
no individual's pay.
\end{dated}

\begin{omfigure}
\begin{tikzpicture}
\begin{axis}[width=10.5cm,height=4.6cm,xmin=50,xmax=800,ymin=0.4,ymax=4.6,ytick=data,
  yticklabels from table={figdata/industry/25-leadership-roles/survey.csv}{label},
  yticklabel style={font=\scriptsize},xticklabel style={font=\footnotesize},xlabel={\small annual wage, \$ thousand},
  grid=major,grid style={gray!20},table/col sep=comma]
\addplot[draw=none,mark=none,error bars/.cd,x dir=both,x explicit,error mark=none,error bar style={omDef,line width=1pt}]
  table[x=p50,y=pos,x error plus expr=\thisrow{p90}-\thisrow{p50},x error minus expr=\thisrow{p50}-\thisrow{p10}]
  {figdata/industry/25-leadership-roles/survey.csv};
\addplot[draw=none,mark=none,error bars/.cd,x dir=both,x explicit,error mark=none,error bar style={omDef!40,line width=7pt}]
  table[x=p50,y=pos,x error plus expr=\thisrow{p75}-\thisrow{p50},x error minus expr=\thisrow{p50}-\thisrow{p25}]
  {figdata/industry/25-leadership-roles/survey.csv};
\addplot[only marks,mark=|,mark size=5pt,line width=1.5pt,omThm] table[x=p50,y=pos]{figdata/industry/25-leadership-roles/survey.csv};
\end{axis}
\end{tikzpicture}

\omcaption{Wages of managers in the securities industry, May 2025 occupational survey: 10th to 90th percentile (thin),
interquartile range (thick), median (mark). Wages exclude bonuses, equity and partners' profit shares. Data:
\texttt{data/industry/oews\_roles.csv}, through \texttt{in\_leader.survey}.}\label{fig:in:leadership-roles:survey}
\end{omfigure}

\section{Accountability regimes}

\begin{definition}[Senior manager function, statement of responsibilities]\label{def:in:leadership-roles:smf}
A \emph{senior manager function}\index{senior manager function} is a role that the UK regulators designate as one whose
holder must be approved before starting and can be held personally accountable for failures in their area. A
\emph{statement of responsibilities}\index{statement of responsibilities} is the single document that sets out what a
senior manager is responsible and accountable for.
\end{definition}

The regime scales with the firm. The FCA's guide lists three functions for limited-scope firms, six for core firms and
seventeen for enhanced firms (\cref{fig:in:leadership-roles:smf}). For a trading firm the core set is typical: a chief
executive (SMF1), executive directors (SMF3) or partners (SMF27), a chair (SMF9), and the required compliance oversight
(SMF16) and money-laundering reporting (SMF17) functions. In the United States, supervisors at a broker-dealer register as
principals; the Series 24 exam ``assesses the competency of an entry-level principal''.

\begin{dated}{2026-09}{Senior managers and principals}\label{dat:in:leadership-roles:smcr}
United Kingdom: FCA guide for solo-regulated firms (July 2019 update): approval before starting for every senior
management function; a \omterm{def:in:leadership-roles:smf}{statement of responsibilities} for every senior manager; 3 functions in the limited-scope tier, 6
in the core tier, 17 in the enhanced tier. The FCA's PS26/6 (first published 21 April 2026) began a reform of the regime:
most phase 1 changes took effect on 24 April 2026, reporting and process changes from 10 July 2026; a second phase needs
legislation. United States: FINRA Series 24, the General Securities Principal exam.
\end{dated}

\begin{omfigure}
\begin{tikzpicture}
\begin{axis}[xbar,width=10.5cm,height=3.8cm,xmin=0,xmax=20,ymin=0.4,ymax=3.6,bar width=11pt,ytick=data,
  yticklabels from table={figdata/industry/25-leadership-roles/smf.csv}{tier},yticklabel style={font=\footnotesize},
  xticklabel style={font=\footnotesize},xlabel={\small senior management functions that can apply},grid=major,
  grid style={gray!20},table/col sep=comma,nodes near coords,nodes near coords style={font=\scriptsize,anchor=west,
  xshift=1pt},point meta=explicit symbolic]
\addplot[fill=omDef,draw=none] table[x=count,y=pos,meta=count]{figdata/industry/25-leadership-roles/smf.csv};
\end{axis}
\end{tikzpicture}

\omcaption{Senior management functions that can apply to FCA solo-regulated firms by tier of the regime, from the FCA's
guide (July 2019 update). A firm allocates those that match its structure and activities. Data:
\texttt{firm.roles.SMF\_TIERS}.}\label{fig:in:leadership-roles:smf}
\end{omfigure}

What the regime changes for the person is accountability: in the guide's words a senior manager ``could be held
accountable if they didn't take reasonable steps to prevent or stop'' a breach in their area, and the \omterm{def:in:leadership-roles:smf}{statement of
responsibilities} says which area is theirs. A quant who becomes a chief risk officer or head of a trading business signs up to that.

\section{Tutorial: making partner}\label{tut:in:leadership-roles:tut}

\begin{tutorial}
\textbf{Goal.} Value a new partner's first five years and compare them with a salaried job; map the regulator's
functions to the firm's roles.
\textbf{End state:} \cref{fig:in:leadership-roles:partner,fig:in:leadership-roles:smf}.
\begin{tutsteps}
\item \textbf{Capital accounts.} Book~16's \texttt{firm.partnership.Partnership} allocates each year's profit by
points, pays the drawn share and retains the rest.
\item \textbf{The path.} \texttt{firm.roles.partner\_path} runs it over simulated profits with a later admission
(\cref{lst:in:leadership-roles:partner}).
\omcode{code/firm/roles/firm_roles.py}{464}{485}{A new partner's draws and capital account over the first years.}\label{lst:in:leadership-roles:partner}
\item \textbf{Break-even.} \texttt{in\_leader.break\_even} bisects the points at which expected draws equal the salaried
alternative, on the same paths.
\item \textbf{Accountability.} \texttt{SMF\_TIERS} and \texttt{smf\_count} give the functions by tier.
\end{tutsteps}
\end{tutorial}

For the example: expected draws \$1.81 million in year one, \$1.76 million in year three after the admissions; capital
account \$2.15 million at the end of year one and \$6.67 million at the end of year five.

\textbf{What to change next.} Let points grow with seniority; repay a departing partner's capital in instalments and
compare staying and leaving in year three; make the firm's profit depend on its capital (Book~16, chapter~2's
treadmill).

\section{Build: leadership cards, accountability and partner economics}\label{bld:in:leadership-roles:roles}

\begin{build}
\bfield{Purpose} Add the leadership roles to the registry, the UK functions by tier, and the new partner's economics.
\bfield{Interface} \texttt{firm.roles}: the cards \texttt{head of desk}, \texttt{chief risk officer}, \texttt{partner},
\texttt{chief executive}; \texttt{SMF\_TIERS}; \texttt{smf\_count(tier)}; \texttt{partner\_path(mean, sd, partners,
points, new\_points, contribution, payout, admit\_year, admit, years, n, rng)} on \texttt{firm.partnership}.
\bfield{Rules} Profit shares by points; a loss reduces capital and pays nothing; admissions dilute from their year; terms
are the caller's, labelled illustrative.
\bfield{Acceptance tests} \texttt{code/firm/roles/tests/}: the tiers have 3, 6 and 17 functions; with a certain profit
the draw is exactly the payout times the share; an admission lowers the share.
\bfield{Stretch} Seniority-weighted points; departures and capital repayment; the partner's certainty equivalent with
the capital at risk (chapter~13).
\end{build}

\begin{omsources}
\item Financial Conduct Authority (2019), The Senior Managers and Certification Regime: Guide for FCA solo-regulated
firms; FCA (2026), PS26/6.
\item FINRA, Series 24; EBA, high earners 2024; BLS occupational survey, May 2025.
\item Book~16, chapters~1, 2, 6 and~12.
\end{omsources}

\section{Exercises}

\begin{exercise}[$\star$]\label{exo:in:leadership-roles:1}
What share of the partnership's points does the new partner hold in year one, and after the admissions?
\end{exercise}

\begin{exercise}[$\star$]\label{exo:in:leadership-roles:2}
How many senior management functions can apply to a core firm, and which two are required functions?
\end{exercise}

\begin{exercise}[$\star$]\label{exo:in:leadership-roles:3}
In a year when the firm makes \$300 million, what does the new partner draw and retain?
\end{exercise}

\begin{exercise}[$\star\star$]\label{exo:in:leadership-roles:4}
Why does the expected draw fall in year three although nothing else changes?
\end{exercise}

\begin{exercise}[$\star\star$]\label{exo:in:leadership-roles:5}
A \omterm{def:in:leadership-roles:head}{head of desk} is offered partnership at 20 points or a raise to \$1.5 million a year. What else must he know?
\end{exercise}

\begin{exercise}[$\star\star$]\label{exo:in:leadership-roles:6}
Why does a partnership keep part of each partner's profit as capital?
\end{exercise}

\begin{exercise}[$\star\star\star$]\label{exo:in:leadership-roles:7}
\emph{Coding.} Raise the payout ratio from 60\% to 80\%. What happens to the expected five-year draws, the year-five
capital account and the probability that capital falls below the contribution?
\end{exercise}

\begin{exercise}[$\star\star\star$]\label{exo:in:leadership-roles:8}
\emph{Find the flaw.} ``Over five years the partner earns \$8.93 million plus \$6.67 million of capital, \$15.6 million
in all: more than twice the \omterm{def:in:leadership-roles:head}{head of desk}.''
\end{exercise}

\section{Problem: Making Partner}

\begin{problem}[{Weekend problem --- making partner}]\label{pb:in:leadership-roles:1}
A \omterm{def:in:leadership-roles:head}{head of desk} at a proprietary trading partnership is offered admission as a partner on the example's terms.

\medskip\noindent\textbf{Part I --- The roles.}
\begin{enumerate}
\item Define the \omterm{def:in:leadership-roles:head}{head of desk}, a \omterm{def:in:leadership-roles:smf}{senior manager function} and a \omterm{def:in:leadership-roles:smf}{statement of responsibilities}.
\item What does a \omterm{def:in:leadership-roles:head}{head of desk} answer for?
\item Which functions do chief officers hold in an enhanced UK firm?
\item What changes when one makes partner?
\item How does a founder's route to the top differ from a manager's?
\end{enumerate}

\medskip\noindent\textbf{Part II --- The model.}
\begin{enumerate}[resume]
\item State the partner's share, draw and retention.
\item State the example's terms.
\item Give the expected five-year draws and their spread.
\item Give the capital account's path and its 5th percentile.
\item How often does the capital fall below the contribution?
\end{enumerate}

\medskip\noindent\textbf{Part III --- Accountability.}
\begin{enumerate}[resume]
\item How many functions apply in each tier?
\item What does a senior manager need before starting?
\item What began in April 2026?
\item What is the US counterpart for supervisors at a broker-dealer?
\item What do the survey and the EBA show about leaders' pay?
\end{enumerate}

\medskip\noindent\textbf{Part IV --- The verdict.}
\begin{enumerate}[resume]
\item State the \emph{named result}: the new partner's expected five-year income and capital at risk, and the
break-even share of profit against staying a salaried \omterm{def:in:leadership-roles:head}{head of desk}.
\item How should the capital account be counted?
\item What does accepting SMF27 add to his responsibilities?
\item What should he negotiate?
\item In two sentences, advise him.
\end{enumerate}
\end{problem}

\section{Interview questions}

\begin{interviewq}[$\star$ \iqroles{trader}]\label{iq:in:leadership-roles:1}
Your desk is 20\% below its revenue budget in October. What do you do?
\end{interviewq}

\begin{interviewq}[$\star$ \iqroles{risk}]\label{iq:in:leadership-roles:2}
As a chief risk officer, how would you set a desk's loss limit?
\end{interviewq}

\begin{interviewq}[$\star\star$ \iqroles{trader, researcher}]\label{iq:in:leadership-roles:3}
Two traders on your desk made the same profit with very different risk. How do you split the bonus pool between them?
\end{interviewq}

\begin{interviewq}[$\star\star$ \iqroles{developer}]\label{iq:in:leadership-roles:4}
As a new chief technology officer, what would you measure in your first month?
\end{interviewq}

\begin{interviewq}[$\star\star$ \iqroles{trader}]\label{iq:in:leadership-roles:5}
When should a \omterm{def:in:leadership-roles:head}{head of desk} stop trading personally?
\end{interviewq}

\begin{interviewq}[$\star\star\star$ \iqroles{trader, risk}]\label{iq:in:leadership-roles:6}
A partner proposes raising the payout ratio to 80\% after a good year. Argue both sides.
\end{interviewq}
